\section{GKT factors inside annular scattering}
\label{sec:gkt-decoration}

The consistent wall structure is determined by the two singleton factors and
the requirement that the path-ordered product about every joint be the
identity.  A
GKT endpoint operation has a different origin: it is the transition between
two systems of open invariants defined by canonical multisections.  The two
constructions use the same Hamiltonian Lie algebra after the $\kappa$-shift,
but their coefficients need not agree.  This section identifies their exact
relation.

\subsection{Labelled critical algebras}

Let $Q=Q_{\GKT}=(\ZZ/4)^2$ be the ordered residue group of
\eqref{eq:GKT-residue-group}.  An object $\mathcal I=(I,\tau)$ consists of a
finite set and a twist map $\tau:I\to Q$; morphisms are twist-preserving
bijections.  Put
\begin{equation}\label{eq:labelled-ring}
 A_{\mathcal I}
 =\QQ[\{u_{i,d}:i\in I, d\geq0\}]
  /\langle u_{i,d}u_{i,d'}:i\in I\rangle.
\end{equation}
Thus two variables belonging to the same marking label have zero product.

Write $\tau(i)=(a_i,b_i)$ with representatives in $\{0,1,2,3\}^2$.  The
quartic GKT critical Lie algebra is spanned over $A_{\mathcal I}$ by
$u_{J,\mathbf d}X_k$, where $J\subset I$ is nonempty,
$u_{J,\mathbf d}=\prod_{i\in J}u_{i,d_i}$, and
\begin{equation}\label{eq:critical-conditions}
 k\equiv\sum_{i\in J}\tau(i)\pmod4,
 \qquad
 k_1+k_2=\sum_{i\in J}\bigl(a_i+b_i+4(d_i-1)\bigr).
\end{equation}
The vector field is
\begin{equation}\label{eq:GKT-generator}
 X_k=x^{k_1}y^{k_2}
 \bigl((k_2+1)x\partial_x-(k_1+1)y\partial_y\bigr).
\end{equation}
The marginal generator $X_{0,0}=x\partial_x-y\partial_y$ is adjoined when
it occurs.

\begin{theorem}[Labelled critical-algebra embedding]
\label{thm:critical-embedding}
For every $\mathcal I$, the assignment
\begin{equation}\label{eq:critical-embedding}
 \jmath_{\mathcal I}:
 u_{J,\mathbf d}X_k
 \longmapsto
 u_{J,\mathbf d}X_{k+\kappa}^{\kappa}
\end{equation}
is a natural coefficient-preserving injective Lie map from the GKT critical
algebra to the Jacobi wall algebra in the distinguished root chamber.  It sends
$X_{0,0}$ to $X_\kappa^\kappa$.  Jacobi transport carries its image to a
strictly path-natural local system on the ordered-residue cover.  The
construction commutes with relabelling, finite restriction,
and deformation truncation.
\end{theorem}

\begin{proof}
In GKT notation,
\[
 [X_k,X_\ell]
 =-\det(k+\kappa,\ell+\kappa)X_{k+\ell}.
\]
Equation~\eqref{eq:shifted-bracket} gives the same formula after applying
\eqref{eq:critical-embedding}.  If two coefficient monomials share a label,
their product vanishes in \eqref{eq:labelled-ring}; if their supports are
disjoint, both conditions in \eqref{eq:critical-conditions} add.  Hence the
assignment is a Lie map, and it is injective on the displayed monomial
basis.  The marginal bracket is preserved because
\[
 [X_{0,0},X_{a,b}]=(a-b)X_{a,b}
 =-\det(\kappa,(a,b)+\kappa)X_{a,b}.
\]
Jacobi path transport is a Lie isomorphism by
Proposition~\ref{prop:Jacobi-transport}.  On the ordered-residue cover the
$Q$-grading is constant; affine transport moves the Hamiltonian exponent and
the $\kappa$-shift together.  The remaining naturality statements are
immediate from the labelled formula.
\end{proof}

In a finite nilpotent quotient, exponential integrates
\eqref{eq:critical-embedding} to an injective group homomorphism.  The GKT
jump theorem identifies an isolated critical zero of a canonical homotopy
with the exponential of its critical-graph generator, while the
boundary-ray theorem of \cite{wuebben2026boundary} gives its unique
increasing-slope factorization, giving the following comparison.

\begin{proposition}[Finite endpoint operations]
\label{prop:endpoint-operations}
Every finite GKT endpoint transition maps to a framed Jacobi automorphism,
and its increasing-slope factors map factor by factor to the shifted rays
$\RR_{\geq0}(k+\kappa)$.  The image depends only on the realizable endpoints,
not on the chosen canonical homotopy.  It is natural under labelled
restriction and path transport.
\end{proposition}

\begin{proof}
Exponentiation and Baker--Campbell--Hausdorff are finite in the chosen
quotient and are functorial for the map of Theorem
\ref{thm:critical-embedding}.  GKT prove that their finite action on
realizable chamber indices is faithful and transitive
\cite[Thm.~4.27 and Cor.~5.14]{grossetal2022open}.  This alone would not
give uniqueness of a transition element.  The enlarged-torsor statement in
Theorem~\ref{thm:companion}, which includes the marginal generator, proves
that the enlarged action is free and transitive.  Its freeness makes
the element carrying one endpoint to the other unique.  The slope assertion
follows because $k$ maps to shifted exponent $k+\kappa$.
\end{proof}

Appendix~\ref{app:gkt-completion} strengthens this formal statement: the
finite endpoint homotopies can be chosen compatibly as the labelled system
grows.  Thus every GKT factor distinguished below has an actual finite
multisection realization.

The ordered-residue cover does not identify annular transport with standard
retwisting of one fixed polynomial critical algebra.  Indeed, the peripheral
residue action completes the singleton twists into
\[
 \{(1,0),(2,3),(3,2),(0,1)\}.
\]
Two distinct labels of twist $(1,0)$ and descendent degree one have a joint
critical block of rank one, with exponent $(2,0)$.  Retwisting both by the
matrix $P_{\mathrm{res}}$ of \eqref{eq:residue-matrix} produces three exponents
$(0,10),(4,6),(8,2)$.  Hence no coefficient-preserving linearization can
fix labels and descendent degrees while retwisting every marking by $P_{\mathrm{res}}$.
Annular transport instead moves the exponent and volume frame and
may produce root-valued Laurent expressions in the target chart.

\subsection{Enumerative and residual factors}

Take two distinct labels with
\[
 t_1=u_{i,1},\quad\tau(i)=(1,0),
 \qquad
 t_2=u_{j,1},\quad\tau(j)=(0,1),
\]
and work modulo $(t_1^2,t_2^2)$.  The two extreme GKT potentials are
\begin{equation}\label{eq:mixed-potentials}
\begin{aligned}
 W^{\min}
 &=x^4+y^4-4t_1xy^4-2t_2y^5+12t_1t_2xy^5,\\
 W^{\max}
 &=x^4+y^4-2t_1x^5-4t_2x^4y+12t_1t_2x^5y.
\end{aligned}
\end{equation}
Their increasing-slope endpoint transition is
\begin{equation}\label{eq:mixed-GKT-transition}
 \exp\!\left(-\frac12t_2X_{0,1}\right)
 \exp\!\left(\frac14t_1t_2X_{1,1}\right)
 \exp\!\left(-\frac12t_1X_{1,0}\right).
\end{equation}
By contrast, the inverse commutator of the two singleton walls has exact
mixed factor
\begin{equation}\label{eq:mixed-total-factor}
 T=\exp\!\left(\frac34t_1t_2X_{1,1}\right).
\end{equation}
Define
\begin{equation}\label{eq:mixed-decoration}
 G_{\mathrm{mix}}=\exp\!\left(\frac14t_1t_2X_{1,1}\right),
 \qquad
 R_{\mathrm{mix}}=\exp\!\left(\frac12t_1t_2X_{1,1}\right).
\end{equation}

\begin{proposition}[First GKT--residual factorization]
\label{prop:first-decoration}
The three factors lie in one abelian shifted-ray group and
\begin{equation}\label{eq:three-quarter-split}
 T=G_{\mathrm{mix}}R_{\mathrm{mix}},
 \qquad \frac34=\frac14+\frac12.
\end{equation}
The first factor on the right is the unique mixed GKT endpoint factor; the
second is the unique residual factor.  Replacing the total wall by the GKT
factor leaves contractible-loop holonomy $R_{\mathrm{mix}}^{\pm1}$.
\end{proposition}

\begin{proof}
The balanced two-label profiles are $(1,5)$ and $(5,1)$.  The GKT chamber
equations give endpoint coefficient $-12$ on each extreme.  After applying
the first singleton factor, the remaining difference in the coefficient of
$xy^5$ is two.  Since
\[
 X_{1,1}(x^4+y^4)=8x^5y-8xy^5,
\]
the mixed endpoint generator has coefficient $2/8=1/4$, proving
\eqref{eq:mixed-GKT-transition}.  Equation~\eqref{eq:first-bracket} gives the
scattering coefficient $3/4$.  The square-zero mixed ray is central, so
subtraction of logarithms gives the residual coefficient $1/2$ exactly.
The total factor is the one which cancels the joint product.
\end{proof}

There is no group projection which fixes the two incoming factors and sends
$T$ to $G_{\mathrm{mix}}$: a homomorphism fixing the inputs fixes their commutator.
Nor can a rank-one character remove $R_{\mathrm{mix}}$, since a character is trivial on the
commutator subgroup.  The GKT factor is therefore a distinguished subfactor
of the total wall, not a quotient wall of a second consistent diagram.

\subsection{All-order factorization}

Let $\mathcal I$ be a finite set of distinct descendent-one labels, each of
twist $\alpha=(1,0)$ or $\beta=(0,1)$, in the square-zero labelled ring.
For $\varnothing\ne J\subseteq\mathcal I$ put
\[
 p_J=|J_\alpha|,
 \qquad q_J=|J_\beta|,
\]
and write
$p_J=r_J+4\ell_J$, $q_J=s_J+4m_J$ with
$0\leq r_J,s_J\leq3$.  The balanced profiles and critical exponents in the
full-support GKT block are
\begin{align}
 b_{J,h}
 &=(r_J+4h,s_J+4(\ell_J+m_J+1-h)),
 &&0\leq h\leq\ell_J+m_J+1,\label{eq:balanced-profiles}\\
 k_{J,h}
 &=(r_J+4h,s_J+4(\ell_J+m_J-h)),
 &&0\leq h\leq\ell_J+m_J.\label{eq:critical-profiles}
\end{align}
Write
\[
 b_J^{\min}=b_{J,0},
 \qquad
 b_J^{\max}=b_{J,\ell_J+m_J+1},
\]
and let $\nu_J^\epsilon$ be the full-support endpoint invariant in the
chamber $\epsilon\in\{\min,\max\}$.  For any exponent $(u,v)$ in this
residue class, set
\begin{equation}\label{eq:quartic-Gamma-weight}
 w_J(u,v)=
 \frac{\Gamma((u+1)/4)\Gamma((v+1)/4)}
      {\Gamma((r_J+1)/4)\Gamma((s_J+1)/4)}.
\end{equation}
This is the quartic Gamma weight in the GKT chamber equation; it is nonzero
for all balanced profiles occurring here.

Let
\begin{equation}\label{eq:GKT-logarithm-J}
 E_J=t_J\sum_hg_{J,h}X_{k_{J,h}}
\end{equation}
be the unique slope-ordered GKT endpoint logarithm.  If $a_{p,q}$ is the
coefficient of $E_{p,q}$ in the unlabelled singleton scattering
factorization, set $a_{p,0}=a_{0,q}=0$ for $p,q>0$ and define the
full-support scattering logarithm
\begin{equation}\label{eq:scattering-logarithm-J}
 S_J=p_J!q_J!a_{p_J,q_J}t_JX_{p_J,q_J}.
\end{equation}
For a shifted ray $\mathfrak r$, let $E_{\mathfrak r}$ and
$S_{\mathfrak r}$ be the sums of the corresponding terms and put
\begin{equation}\label{eq:residual-ray}
 R_{\mathfrak r}=S_{\mathfrak r}-E_{\mathfrak r}.
\end{equation}

\begin{theorem}[Canonical all-order GKT factorization]
\label{thm:all-order-decoration}
The endpoint recursions determine every coefficient in
\eqref{eq:GKT-logarithm-J} canonically, and
\begin{equation}\label{eq:decorated-word}
 \overrightarrow{\prod_{\mathfrak r}}\exp(S_{\mathfrak r})
 =\overrightarrow{\prod_{\mathfrak r}}
  \bigl(\exp(E_{\mathfrak r})\exp(R_{\mathfrak r})\bigr).
\end{equation}
The equality is raywise; no factor is commuted past a factor on another
slope.  It is natural under relabelling, deletion of labels, restriction of
finite labelled systems closed under the dependencies of
Appendix~\ref{app:gkt-completion}, reduction of deformation order, and framed
Jacobi transport.  It adds no new geometric support to the annular diagram.
\end{theorem}

\begin{proof}
The GKT dimension formula gives
\eqref{eq:balanced-profiles}--\eqref{eq:critical-profiles}.  For an endpoint
$\epsilon\in\{\min,\max\}$, the full-support invariant is determined by the
partition recursion
\begin{equation}\label{eq:endpoint-recursion}
 w_J(b_J^\epsilon)\nu_J^\epsilon
 +\sum_{\substack{\pi\in\operatorname{Part}(J)\\|\pi|\geq2}}
 w_J\!\left(\sum_{C\in\pi}b_C^\epsilon\right)
 \prod_{C\in\pi}\nu_C^\epsilon
 =\begin{cases}-1,&|J|=1,\\0,&|J|\geq2.
 \end{cases}
\end{equation}
The coefficient of $\nu_J^\epsilon$ is a nonzero Gamma weight, so induction
on $|J|$ determines the endpoints.  The critical generators form the
weighted incidence matrix of the path of balanced profiles.  Indeed,
\begin{equation}\label{eq:critical-incidence}
 X_{u,v}(x^4+y^4)
 =4(v+1)x^{u+4}y^v-4(u+1)x^uy^{v+4}.
\end{equation}
For $(u,v)=k_{J,h}$ the two monomials in
\eqref{eq:critical-incidence} are precisely the adjacent vertices
$b_{J,h+1}$ and $b_{J,h}$.  The resulting path incidence matrix has full
column rank, and its one left-kernel relation is the chamber equation.
Hence the coefficients $g_{J,h}$ exist and are unique.

For the scattering term, substitute
\[
 t_\alpha=\sum_{\tau(i)=\alpha}t_i,
 \qquad t_\beta=\sum_{\tau(i)=\beta}t_i
\]
in the unlabelled factorization.  The coefficient of $t_J$ in
$t_\alpha^{p_J}t_\beta^{q_J}$ is $p_J!q_J!$, which gives
\eqref{eq:scattering-logarithm-J}.  Both constructions are triangular in
proper subsets, proving naturality under restriction.  Generators on one
shifted ray commute, so \eqref{eq:residual-ray} exponentiates and yields
\eqref{eq:decorated-word}.  A nonzero $E_{\mathfrak r}$ with
$S_{\mathfrak r}=0$ is recorded as the identity pair
$\exp(E_{\mathfrak r})\exp(-E_{\mathfrak r})$, not as a geometric wall.
\end{proof}

\subsection{Compatible geometric realization}

The preceding factorization has a geometric realization.  The relative
extension arguments in Appendix~\ref{app:gkt-completion} give the following
geometric realization theorem.

For the theorem, an \emph{exact GKT datum} is a graded open $(4,4)$-graph
together with its descendent vector.  A \emph{finite partial system} records
a finite twist-labelled set, a finite set of labelled monomials closed under
divisors, and finitely many exact data closed under the connected-component
and boundary dependencies of the GKT assembling relation.  The precise
definition is given in Appendix~\ref{app:gkt-completion}.

\begin{theorem}[Compatible completed GKT realization]
\label{thm:compatible-GKT-realization}
There are nested finite partial systems
$\mathbb D_1\subset\mathbb D_2\subset\cdots$ exhausting all exact GKT data,
compatible canonical families realizing the two extremal chamber indices,
and a componentwise smooth GKT homotopy between them such that:
\begin{enumerate}
\item the countably many ray factors occur in pairwise disjoint open time
intervals, ordered by increasing slope;
\item each fixed component is nonconstant on only finitely many such
intervals and has constant endpoint collars;
\item every finite divisor-closed quotient has precisely its prescribed
finite collection of nonzero ray jumps; and
\item the finite restrictions form an inverse limit in the category of
sets.
\end{enumerate}
Simplicity is asserted for the finite elementary homotopies, not for one
global assignment of jump times to all label subsets.
\end{theorem}

Thus every GKT factor in Theorem~\ref{thm:all-order-decoration} is the image
of an actual finite restriction of a compatible homotopy.  The residual
factors are not enumerative jumps; they remain comparison automorphisms.

\subsection{Residual coherence}

The residual factors can be retained coherently in multiplication.  Fix a
finite labelled system and two framed input paths.  Let $U_i$ be the
total consistent transport, $V_i$ its GKT subword, and
$Q_i=U_iV_i^{-1}$ the residual transport.  In a finite nilpotent quotient
define the midpoint
\begin{equation}\label{eq:residual-midpoint}
 H=Q_1(Q_1^{-1}Q_2)^{1/2},
 \qquad C_i=HQ_i^{-1}.
\end{equation}
Rational powers are the finite polynomials $g^c=\exp(c\log g)$.

\begin{lemma}\label{prop:residual-midpoint}
The midpoint is symmetric in $Q_1,Q_2$ and
\begin{equation}\label{eq:corrected-product}
 C_iU_i=HV_i,
 \qquad
 \mu(C_1U_1f_1,C_2U_2f_2)=H\mu(V_1f_1,V_2f_2).
\end{equation}
It is natural under finite restriction and bi-equivariant under common
terminal total and GKT transports.
\end{lemma}

\begin{proof}
With $A=Q_1^{-1}Q_2$,
$Q_2(Q_2^{-1}Q_1)^{1/2}=Q_1A(A^{-1})^{1/2}=Q_1A^{1/2}$, proving symmetry.
The first identity in \eqref{eq:corrected-product} is immediate; the second
uses that $H$ is an algebra automorphism.  Conjugation-equivariance of
rational powers gives the claimed bi-equivariance.
\end{proof}

For positive degrees $r,s$, set
\begin{equation}\label{eq:weighted-interpolation}
 M_{r,s}(P,Q)=P(P^{-1}Q)^{s/(r+s)}.
\end{equation}
For a rooted binary product tree, recursively define a residual interpolant,
with weights equal to the sums of the leaf degrees.

\begin{proposition}[Degree-weighted coherence]
\label{thm:residual-coherence}
If $H_\tau$ is the interpolant associated with a product tree $\tau$, its
corrected product is
\[
 \mathsf P_\tau=H_\tau\prod_iV_if_i.
\]
For two trees $\tau,\sigma$, the automorphism
$\mathcal A_{\tau\leftarrow\sigma}=H_\tau H_\sigma^{-1}$ carries one
product to the other and satisfies
\[
 \mathcal A_{\tau\leftarrow\sigma}
 \mathcal A_{\sigma\leftarrow\rho}
 =\mathcal A_{\tau\leftarrow\rho}.
\]
All associahedral coherence identities therefore hold.  For three
equal-degree residuals $e^x,e^y,e^z$ in the free labelled square-zero group,
\begin{equation}\label{eq:first-associator}
 \log\mathcal A_{L\leftarrow R}=\frac1{72}[[x,z],y].
\end{equation}
The comparison is coherent but not generally strict monoidal.
\end{proposition}

\begin{proof}
At an internal vertex, the corrections carry both child products to the
same interpolant $H_\tau$.  Since this interpolant is an algebra automorphism, induction
gives the displayed form of $\mathsf P_\tau$.  Ratios of interpolants carry one
tree product to another, and cancellation proves the cocycle identity.
The weighted interpolation is symmetric under exchange of the two weighted
children and is equivariant under the left--right bitorsor action.
Baker--Campbell--Hausdorff expansion for three equal weights gives the common
linear term $(x+y+z)/3$, no quadratic difference, and cubic difference
\[
 \frac1{72}(xzy-zxy-yxz+yzx)
 =\frac1{72}[[x,z],y],
\]
which is \eqref{eq:first-associator}.  By comparison, unweighted iteration
assigns leaf weights $(1/4,1/4,1/2)$ and $(1/2,1/4,1/4)$ in the two
parenthesizations already in the abelian quotient.
\end{proof}

Changes of framed input paths act on the individual legs as well as on the
output interpolant.  The precise naturality statement is the following.

Let $\Gamma=(\gamma_i)$ and $\Gamma'=(\gamma_i')$ be two tuples of framed
input paths with the same endpoints.  For a fixed product tree,
write $V_i^\Gamma,V_i^{\Gamma'}$ for the corresponding GKT transports and
$H_\Gamma,H_{\Gamma'}$ for their residual interpolants.  Put
\begin{equation}\label{eq:corridor-maps}
 W_i^{\Gamma'\leftarrow\Gamma}
   =V_i^{\Gamma'}(V_i^\Gamma)^{-1},
 \qquad
 B^{\Gamma'\leftarrow\Gamma}
   =H_{\Gamma'}H_\Gamma^{-1}.
\end{equation}

\begin{lemma}[Change of framed paths]
\label{prop:corridor-groupoid}
For $a_i=V_i^\Gamma f_i$ one has the commutative comparison square
\begin{equation}\label{eq:corridor-square}
 H_{\Gamma'}\prod_iV_i^{\Gamma'}f_i
 =B^{\Gamma'\leftarrow\Gamma}H_\Gamma
   \prod_i\bigl(W_i^{\Gamma'\leftarrow\Gamma}a_i\bigr).
\end{equation}
The leg maps and output maps compose under successive changes of framed paths, and
the resulting squares commute with the tree-transition cocycle of
Proposition~\ref{thm:residual-coherence}.  Thus path and tree changes form a
coherent factorization groupoid.

There is no output-only descent for arbitrary independent changes of the input paths.
In particular, if one leg changes by a nontrivial unital automorphism $W$
and another leg is fixed, no output automorphism $Z$ can satisfy
$Z(ab)=W(a)b$ for all $a,b$.
\end{lemma}

\begin{proof}
The identities
$V_i^{\Gamma'}f_i=W_i^{\Gamma'\leftarrow\Gamma}a_i$ and
$B^{\Gamma'\leftarrow\Gamma}H_\Gamma=H_{\Gamma'}$ give
\eqref{eq:corridor-square}.  For a third tuple $\Gamma''$, cancellation
gives
\[
 W_i^{\Gamma''\leftarrow\Gamma'}
 W_i^{\Gamma'\leftarrow\Gamma}
 =W_i^{\Gamma''\leftarrow\Gamma},
 \qquad
 B^{\Gamma''\leftarrow\Gamma'}
 B^{\Gamma'\leftarrow\Gamma}
 =B^{\Gamma''\leftarrow\Gamma}.
\]
Tree transitions are likewise ratios of the corresponding centers, so the
mixed squares commute by the same cancellation.  Finally, if
$Z(ab)=W(a)b$, then $b=1$ gives $Z=W$, whereas $a=1$ gives
$Z=\mathrm{id}$.  Hence an output-only comparison exists only when the leg
change is trivial.
\end{proof}

The result of this section is therefore a coherent factorization comparison.
The total factors remain the walls of the consistent diagram.  The GKT
endpoint factors are marked inside them, and the residual automorphisms are
retained as comparison data.  No equality between these wall coefficients
and commutative theta structure constants has been asserted.
